\section{Method}
\subsection{Simulation Framework}
\subsubsection{Environment}

We model the problem as a partially observable stochastic game defined as a tuple
$\langle N, S, A, T, R, \Omega, O \rangle$.
$N = 6$ agents interact over a shared depletable common-pool resource.
The state space $S$ is defined by the current resource level $r_t \geq 0$,
and each agent $i \in \{1, \ldots, N\}$ selects an action
$a_i^t \in A = \{0, 1\}$ (0: abstain, 1: exploit) at each timestep $t$,
yielding joint action $\mathbf{a}_t = (a_1^t, \ldots, a_N^t)$.

Agents have partial observability: each agent may receive, as observation
$o_i \in \Omega$, an affective signal broadcast by other agents
(Section~\ref{sec:emotions}), or no observation at all.

\subsection{Resource Dynamics}

For each agent $i$ with $a_i^t = 1$,
success occurs with probability proportional to resource availability:
\begin{equation}
    P(\text{success}_i \mid r_t) = \frac{r_t}{R_0}
\end{equation}
where $R_0$ is the initial resource level.

Let $C_t$ denote the total number of successful exploitations at time $t$. The resource evolves as follows:
\begin{equation}
    r_{t+1} =  r_t - C_t
\end{equation}
This formulation induces a natural commons dilemma: individual exploitation pressure
degrades collective future access.

\subsection{Affective Signaling}
\label{sec:emotions}


We introduce a lightweight inter-agent communication mechanism based on
scalar \emph{affective signals}, grounded in appraisal theory
\cite{moorsPowerGoalDirectedProcesses2017}: each agent broadcasts a valence signal derived from its
recent consumption rate $r_i^{\text{cons}} \in [0,1]$.

$r_i^{\text{cons}}$ is the proportion of the last $m = 10$
timesteps in which agent $i$ successfully exploited the resource. Personal satisfaction is
then defined as:
\begin{equation}
    r_i^{\text{sat}} = \beta \cdot \mathbf{1}[\text{last\_meal}_i]
                     + (1 - \beta) \cdot r_i^{\text{cons}}, \quad \beta = 0.5
\end{equation}
blending the most recent consumption outcome with the consumption rate.

The valence signal broadcast by each agent, $e_i \in [-1, 1]$, is defined as:
\begin{equation}
    e_i = 2\,r_i^{\text{cons}} - 1
\end{equation}

Each agent $i$ receives, as observation $o_i$, the mean affective signal
of all other agents:
\begin{equation}
    o_i = \frac{1}{N-1}\sum_{j \neq i} e_j
\end{equation}

\subsection{Reward Function}
\label{sec:reward}

The reward function implements \emph{empathy-driven reward shaping}
\cite{lynneTheoreticalFoundationEmpathy2016}. For a successful exploitation, each agent $i$ receives:
\begin{equation}
    r_i = r_i^{\text{sat}} \cdot (1 - \alpha_e) \;+\; \alpha_e \cdot o_i
\end{equation}
where:
\begin{itemize}
    \item $r_i^{\text{sat}} \in [0, 1]$: agent's personal satisfaction
          (see Section~\ref{sec:emotions}),
    \item $o_i \in [-1, 1]$: mean affective signal received from other agents
          (see Section~\ref{sec:emotions}),
    \item $\alpha_e \in [0, 1]$: empathic weight, interpolating between
          purely self-interested ($\alpha_e = 0$) and purely
          other-regarding ($\alpha_e = 1$) behaviour.
\end{itemize}

\subsection{Learning Algorithm}

All agents are trained via Deep Q-Network (DQN) \cite{mnih2015}, each taking the scalar observation $o_i$ as sole input. All agents share the same network architecture and global hyperparameters (see Appendix~\ref{app:hyperparams}), but each agent maintains its own Q-network weights updated from its individual experience. Agents follow an $\varepsilon$-greedy exploration policy, and the Q-network is updated by minimising the mean squared Bellman error over sampled experience batches.

\subsection{Study 1: Factorial Manipulation of Observation and Empathic Reward}

\subsubsection{Framework}

Study 1 adopts a $2\times2$ factorial design that crosses two independent
dimensions: \emph{observation} (whether agents receive affective signals
$o_i$ or a zero vector) and \emph{empathic weight}
($\alpha_e \in \{0,\, 0.5\}$). This yields four homogeneous conditions,
summarised in Table~\ref{tab:conditions}.

\begin{table}[h]
\centering
\begin{tabular}{llcc}
\hline
\textbf{Label} & \textbf{Description} & \textbf{Observation} & $\boldsymbol{\alpha_e}$ \\
\hline
Blind / Non-empathic & Baseline             & No  & $0.0$ \\
Blind / Empathic     & Reward only          & No  & $0.5$ \\
Sees / Ignores       & Observation only     & Yes & $0.0$ \\
Full Empathy         & Observation + Reward & Yes & $0.5$ \\
\hline
\end{tabular}
\caption{Conditions in Study 1.}
\label{tab:conditions}
\end{table}

Common simulation parameters are listed in Appendix~\ref{app:simparams}.

\subsubsection{Dependent Variables and Statistical Analysis}

Two dependent variables are measured per episode:
\begin{itemize}
    \item \textbf{Episode length}: the number of steps before the resource
          is fully depleted, used as a proxy for collective sustainability.
          Longer episodes indicate that agents collectively conserved the
          resource more effectively.
    \item \textbf{Fairness}: the Gini coefficient of resource access across
          agents within an episode, measuring inequality in the distribution
          of successful exploitations. Lower values indicate more equitable
          access.
\end{itemize}

A $2 \times 2$ between-subjects ANOVA is run on each dependent variable,
with two factors:
\begin{itemize}
    \item \textbf{Observation} (absent vs.\ present): whether agents receive
          the mean affective signal $o_i$ or a zero vector.
    \item \textbf{Empathic weight} ($\alpha_e = 0$ vs.\ $\alpha_e = 0.5$):
          whether the reward function includes a social component.
\end{itemize}
We test the main effect of each factor and their interaction.

\subsection{Study 2: Graded Empathic Weight}

\subsubsection{Framework}

Study 2 holds observation fixed --- all agents receive the mean affective
signal of their peers --- and varies $\alpha_e$ across five levels:
$\{0.0,\, 0.25,\, 0.5,\, 0.75,\, 0.99\}$. All other parameters are
identical to Study 1, with four independent replications per $\alpha_e$ level.
This characterises the dose-response relationship between empathic weight and
collective resource management.

\subsubsection{Dependent Variables and Statistical Analysis}

The same two dependent variables as Study 1 are recorded (episode length
and Gini coefficient of resource access). Because $\alpha_e$ is treated
as a continuous ordinal variable rather than a categorical factor, the
monotonic relationship between $\alpha_e$ and each outcome is assessed
using Spearman's rank correlation $r_s$.
